% संपूर्ण मराठी ग्रंथातील प्रकरण 32: विन्यासमीमांसा आणि चिन्हार्थमीमांसा
% हा संपादनयोग्य स्रोतखंड आहे. संकलनासाठी संपूर्ण स्रोतसंग्रहातील
% अक्षररूपे, पूर्वभाग, चिन्हव्याख्या आणि आकृती-स्रोत आवश्यक आहेत.
% मूळ: Open Logic Project; मजकूर आणि रूपांतर CC BY 4.0.
% यंत्रानुवाद, दुरुस्ती आणि तपासणी: OpenAI Codex —
% GPT-5.6 Sol आणि GPT-6 Sol, दोन्ही Ultra effort.
% दुरुस्ती-पुनरावलोकन: GPT-6.1 Sol, Ultra effort.
\chapter{विन्यासमीमांसा आणि चिन्हार्थमीमांसा}\label{sol:syn:chap}
\begin{editorial}
हा द्वितीय-क्रम
तर्कशास्त्रावरील
भागाचा आरंभ आहे.
\end{editorial}
\begin{editorial}
आतापर्यंत येथे
द्वितीय-क्रम
तर्कशास्त्राची
मूलभूत विन्यासमीमांसा
आणि चिन्हार्थमीमांसा
घेतली आहे.
स्वतंत्र प्रकरण
म्हणून हे फारच
लहान आहे.
द्वितीय-क्रम
चलांसाठीचे आदेशन
घेतल्याशिवाय
द्वितीय-क्रम
तर्कशास्त्राच्या
निष्पत्ती प्रणालींबद्दल
बोलता येणार नाही;
आणि तेथे काही
सूक्ष्म प्रश्नही
उद्भवतात.
\end{editorial}









\section{परिचय}\label{sol:syn:int:sec}

प्रथम-क्रम तर्कशास्त्रात,
दिलेल्या भाषेतील
अतार्किक चिन्हे,
म्हणजे तिचे
स्थिरांक,
फलने
आणि विधेये,
तार्किक चिन्हांबरोबर
जोडून आपण
प्रथम-क्रम
संरचना
विषयी विधाने
मांडतो.
हे पूर्ति
या संकल्पनेने
केले जाते.
ती संरचना~
$\Struct{M}$,
चर-मूल्यांकन~$s$
आणि सूत्र~$\varphi$
यांना जोडते:
$\Sat{M}{\varphi}[s]$
खरे असते, जर
आणि फक्त जर
त्यातील
स्थिरांक,
फलने
आणि विधेये
यांचा अर्थ
$\Struct{M}$ नुसार
आणि त्यातील मुक्त
चलांचा अर्थ~$s$
नुसार घेतल्यावर
$\varphi$ जे व्यक्त
करते ते सत्य
असते.
एकरूपता~$\eq$
चा अर्थ
$\Sat{M}{\varphi}[s]$
च्या व्याख्येतच
दिलेला असतो;
$\lforall$ आणि
$\lexists$ यांचाही
अर्थ तसाच
दिलेला असतो.
यांपैकी पहिले
चिन्ह रचनेच्या
प्रांत~$\Domain{M}$
वरील एकरूपता
संबंध म्हणूनच
नेहमी अर्थ लावले
जाते; संख्यकांचा
आवाका नेहमी
संपूर्ण प्रांत
असतो.
पण महत्त्वाचे
म्हणजे संख्यकीकरण
फक्त प्रांत
मधील घटकांवर
करता येते;
म्हणून संख्यकानंतर
फक्त वस्तू-चले
येऊ शकतात.

द्वितीय-क्रम
तर्कशास्त्रात
भाषा आणि पूर्तीची
व्याख्या या दोन्ही
विस्तारून त्यांत
मुक्त व बद्ध
फलन-चल, विधेय-चल
आणि त्यांवरील
संख्यकीकरण
समाविष्ट केले
जाते.
वस्तू-चलांचा
स्थिरांक शी
जसा संबंध असतो,
तसाच या चलांचा
फलने
आणि विधेये
शी असतो.
द्वितीय-क्रम
तर्कशास्त्रातील
पदे व सूत्रे
घडवताना त्यांची
तशीच भूमिका
असते आणि
त्यांवरील
संख्यकीकरणही
तत्सम रीतीने
हाताळले जाते.
\emph{मानक}
चिन्हार्थमीमांसेत
द्वितीय-क्रम
संख्यकांचा
आवाका योग्य
प्रकारच्या सर्व
संभाव्य वस्तूंवर
असतो (फलन-चलांसाठी
$\Domain{M}$ वरून~$\Domain{M}$
मध्ये जाणारी
$n$-स्थानी फलने,
आणि विधेय-चलांसाठी
$n$-स्थानी संबंध).
उदाहरणार्थ,
$\lforall[\Obj{v_0}][(\Obj{P^1_0}(\Obj{v_0}) \lor \lnot
  \Obj{P^1_0}(\Obj{v_0}))]$
हे प्रथम-क्रम
आणि द्वितीय-क्रम
या दोन्ही
तर्कशास्त्रांतील
सूत्र आहे;
पण दुसऱ्यामध्ये
$\lforall[\Obj{V^1_0}][\lforall[\Obj{v_0}][(\Obj{V^1_0}(\Obj{v_0}) \lor
    \lnot \Obj{V^1_0}(\Obj{v_0}))]]$
आणि
$\lexists[\Obj{V^1_0}][\lforall[\Obj{v_0}][(\Obj{V^1_0}(\Obj{v_0}) \lor
    \lnot \Obj{V^1_0}(\Obj{v_0}))]]$
यांचाही विचार
करता येतो.
त्यांत मुक्त
चले नसल्याने
ती द्वितीय-क्रम
तर्कशास्त्रातील
वाक्ये
आहेत.
येथे $\Obj{V^1_0}$
हे द्वितीय-क्रम
$1$-स्थानी
विधेय-चल आहे.
$\Obj{V^1_0}$
चे अनुमत
अर्थनिर्धारण
$1$-स्थानी
विधेय
असलेल्या
$\Obj{P^1_0}$
ला देता येणाऱ्या
अर्थनिर्धारणासारखेच
असते, म्हणजे
ते~$\Domain{M}$
चे उपसंच असतात.
त्यांवर संख्यकीकरण
केल्यास
$\lforall[\Obj{v_0}][(\Obj{V^1_0}(\Obj v_0) \lor \lnot
  \Obj{V^1_0}(v_0))]$
हे $\Obj{V^1_0}$
ला~$\Domain{M}$
चा उपसंच मूल्य
म्हणून देण्याच्या
प्रत्येक पद्धतीत
किंवा किमान
एका पद्धतीत
खरे आहे असे
म्हटले जाते.
प्रत्येक संचात
दिलेला घटक
असतो किंवा
नसतो; म्हणून
दोन्ही वाक्ये
कोणत्याही
संरचना
मध्ये सत्य असतात.











\section{पदे आणि सूत्र}\label{sol:syn:frm:sec}

प्रथम-क्रम
तर्कशास्त्राप्रमाणेच
द्वितीय-क्रम
तर्कशास्त्रातील
अभिव्यक्तींची
रचना मूलभूत
शब्दसंग्रहापासून
होते. त्यात
\emph{चले},
\emph{स्थिरांक},
\emph{विधेये}
आणि कधीकधी
\emph{फलने}
असतात.
या चिन्हांपासून,
तार्किक संयोजके,
संख्यक आणि कंस
व स्वल्पविराम
यांसारखी
विरामचिन्हे
यांच्या साहाय्याने,
\emph{पदे} आणि
\emph{सूत्रे}
तयार होतात.
फरक असा की
वस्तूंसाठीच्या
चलांव्यतिरिक्त
द्वितीय-क्रम
तर्कशास्त्रात
संबंधांसाठी
आणि फलनांसाठीही
चले असतात;
त्यांवर
संख्यकीकरण
करण्याचीही
मुभा असते.
म्हणून द्वितीय-क्रम
तर्कशास्त्राची
तार्किक चिन्हे
म्हणजे प्रथम-क्रम
तर्कशास्त्राची
चिन्हे आणि
यांखेरीज पुढील:

\begin{enumerate}
\item प्रत्येक
  स्थानसंख्या~$n$
  साठी द्वितीय-क्रम
  संबंध-चले
  चा गणनीय अनंत
  संच:
  $\Obj V_0^n$, $\Obj V_1^n$, $\Obj V_2^n$, \dots
\item द्वितीय-क्रम
  फलन-चले
  चा गणनीय अनंत
  संच:
  $\Obj u_0^n$, $\Obj u_1^n$, $\Obj u_2^n$, \dots
\end{enumerate}

प्रथम-क्रम
चल~$\Obj v_i$
यांसाठी आपण
$x$, $y$, $z$
ही अधिचले
वापरतो; त्याचप्रमाणे
$\Obj V_i^n$
यांसाठी
$X$, $Y$, $Z$
इत्यादी, आणि
$\Obj u_i^n$
यांसाठी
$u$, $v$
इत्यादी अधिचले
वापरू.

\begin{explain}
द्वितीय-क्रम
भाषेतील अतार्किक
चिन्हे प्रथम-क्रम
भाषेप्रमाणेच
निर्दिष्ट होतात:
तिची स्थिरांक,
फलने
आणि विधेये
यांची यादी देऊन.

प्रथम-क्रम
तर्कशास्त्रात
एकरूपता~$\eq$
सहसा समाविष्ट
असते.
प्रथम-क्रम
तर्कशास्त्रात
भाषा~$\Lang{L}$
मधील अतार्किक
चिन्हे काही
लक्षणीय गोष्ट
व्यक्त करण्यासाठी
अत्यावश्यक असतात.
अतार्किक चिन्हे
न वापरणारी
वाक्ये
अर्थात असतात;
पण फक्त~$\eq$
असल्यास काही
लक्षणीय गोष्ट
सांगणे कठीण
असते.
द्वितीय-क्रम
तर्कशास्त्रात
आपल्याकडे संबंध-
आणि फलन-चलांचा
अमर्याद पुरवठा
असल्यामुळे,
अतार्किक चिन्हांचा
विशेष पुरवठा
नसतानाही
प्रथम-क्रम
भाषेत जे काही
सांगता येते
ते सांगता येते.
\end{explain}

\begin{defn}[द्वितीय-क्रम पदे]
भाषा~$\Lang L$
मधील \emph{द्वितीय-क्रम
पदांचा} संच
$\TrmSOL[L]$
हा \ref{fol:syn:frm:defn:terms}
मध्ये पुढील
खंड वाढवून
परिभाषित होतो:
\begin{enumerate}
\item $u$ हे
  $n$-स्थानी
  फलन-चल
  आणि $t_1$,
  \dots, $t_n$
  ही पदे असतील,
  तर $\Atom{u}{t_1, \ldots, t_n}$
  हे पद असते.
\end{enumerate}
\end{defn}

\begin{explain}
म्हणून द्वितीय-क्रम
पद प्रथम-क्रम
पदासारखेच दिसते;
फक्त प्रथम-क्रम
पदात जिथे
फलन~$\Obj{f^n_i}$
असते, तिथे
द्वितीय-क्रम
पदात त्याच्या
जागी फलन-चल~$\Obj{u^n_i}$
असू शकते.
\end{explain}

\begin{defn}[द्वितीय-क्रम सूत्र]
भाषा~$\Lang L$
मधील \emph{द्वितीय-क्रम
सूत्रे} चा
संच~$\FrmSOL[L]$
हा \ref{fol:syn:frm:defn:formulas}
मध्ये पुढील
खंड वाढवून
परिभाषित होतो:
\begin{enumerate}
\item $X$ हे
  $n$-स्थानी
  विधेय-चल
  आणि $t_1$,
  \dots, $t_n$
  ही~$\Lang L$
मधील द्वितीय-क्रम
पदे असतील, तर
  $\Atom{X}{t_1,\ldots, t_n}$
  हे अणुरूप
  सूत्र
  असते.

\item $\varphi$ हे सूत्र आणि $u$ हे फलन-चल असेल,
  तर $\lforall[u][\varphi]$ हे सूत्र असते.

\item $\varphi$ हे सूत्र आणि $X$ हे विधेय-चल असेल,
  तर $\lforall[X][\varphi]$ हे सूत्र असते.

\item $\varphi$ हे सूत्र आणि $u$ हे फलन-चल असेल,
  तर $\lexists[u][\varphi]$ हे सूत्र असते.

\item $\varphi$ हे सूत्र आणि $X$ हे विधेय-चल असेल,
  तर $\lexists[X][\varphi]$ हे सूत्र असते.
\end{enumerate}
\end{defn}












\section{पूर्ति}\label{sol:syn:sat:sec}

\begin{explain}
द्वितीय-क्रम
सूत्रे साठी
पूर्ति-संबंध
$\Sat{M}{\varphi}[s]$
परिभाषित करताना
द्वितीय-क्रम
चले
समाविष्ट करण्यासाठी
व्याख्या वाढवाव्या
लागतात.
संरचना
ही संकल्पना
द्वितीय-क्रम
आणि प्रथम-क्रम
तर्कशास्त्रात
सारखीच असते.
फरक फक्त
चर-मूल्यांकन~$s$
बाबत असतो:
आता त्याने
प्रथम-क्रम
चले
बरोबरच
द्वितीय-क्रम
चले
यांनाही मूल्ये
द्यावी लागतात.
\end{explain}

\begin{defn}[चर-मूल्यांकन]
संरचना~$\Struct{M}$
साठीचे
\emph{चर-मूल्यांकन}~$s$
हे असे फलन
आहे जे पुढीलपैकी
प्रत्येकाला
दिल्याप्रमाणे
प्रतिचित्रित करते:
\begin{enumerate}
\item वस्तू-चल~$\Obj{v_i}$
  याला~$\Domain M$
  मधील घटकाकडे,
  म्हणजे
  $s(\Obj{v_i}) \in \Domain{M}$;
\item $n$-स्थानी
  संबंध-चल~$\Obj{V_i^n}$
  याला~$\Domain{M}$
  वरील $n$-स्थानी
  संबंधाकडे,
  म्हणजे
  $s(\Obj{V_i^n}) \subseteq \Domain{M}^n$;
\item $n$-स्थानी
  फलन-चल~$\Obj{u_i^n}$
  याला
  $\Domain{M}$
  वरून $\Domain{M}$
  मध्ये जाणाऱ्या
  $n$-स्थानी
  फलनाकडे,
  म्हणजे
  $s(\Obj{u_i^n})\colon \Domain{M}^n \to \Domain{M}$;
\end{enumerate}
\end{defn}

\begin{explain}
संरचना
प्रत्येक स्थिरांक
आणि फलन
ला मूल्य
देते, तर द्वितीय-क्रम
चर-मूल्यांकन
प्रत्येक वस्तू-चल
आणि फलन-चल
यांना अनुक्रमे
वस्तू आणि फलने
देते. या दोन्हींच्या
साहाय्याने
प्रत्येक पदाला
मूल्य देता येते.
\end{explain}

\begin{defn}[पदाचे मूल्य]
$t$ हे भाषा~$\Lang L$
मधील पद,
$\Struct M$ ही~$\Lang L$
साठी संरचना,
आणि $s$ हे~$\Struct M$
साठी चल
मूल्यांकन असेल,
तर \emph{मूल्य}~
$\Value{t}{M}[s]$
हे प्रथम-क्रम
पदांप्रमाणेच,
पुढील अतिरिक्त
खंडासह परिभाषित
होते:
\begin{quote}
\indcase{t}{\Atom{u}{t_1, \ldots, t_n}}{
\[\Value{\indfrm}{M}[s] = s(u)(\Value{t_1}{M}[s], \ldots,
\Value{t_n}{M}[s]).
\]}
\end{quote}
\end{defn}

\begin{defn}[$x$-पर्याय]
$s$ हे संरचना~$\Struct M$
साठी चल
मूल्यांकन असेल,
तर~$\Struct M$
साठीचे असे
कोणतेही चल
मूल्यांकन $s'$
जे $s$ पेक्षा
जास्तीत जास्त
$x$ ला दिलेल्या
मूल्यातच भिन्न
असेल, त्याला
$s$ चा
\emph{$x$-पर्याय}
म्हणतात.
$s'$ हा $s$
चा $x$-पर्याय
असेल, तर
$\varAssign{s'}{s}{x}$
असे लिहितो.
(द्वितीय-क्रम
चल $X$ किंवा
$u$ यांसाठीही
तसेच.)
\end{defn}

\begin{defn}

  $s$ हे संरचना~$\Struct M$
  साठी चल
  मूल्यांकन
  आणि $m \in \Domain{M}$
  असेल, तर
  मूल्यांकन~$\Subst{s}{m}{x}$
  हे पुढीलप्रमाणे
  परिभाषित केलेले
  चल
  मूल्यांकन आहे:
  \[\Subst{s}{m}{x}(y) = \begin{cases}
    m & \text{जर } y \ident x\\
    s(y) & \text{अन्यथा},
  \end{cases}\]
  $X$ हे $n$-स्थानी
  संबंध-चल
  आणि $M \subseteq
  \Domain{M}^n$
  असेल, तर
  $\Subst{s}{M}{X}$
  हे पुढीलप्रमाणे
  परिभाषित केलेले
  चल
  मूल्यांकन आहे:
  \[\Subst{s}{M}{X}(y) = \begin{cases}
    M & \text{जर } y \ident X\\
    s(y) & \text{अन्यथा}.
  \end{cases}\]
  $u$ हे $n$-स्थानी
  फलन-चल
  आणि $f\colon \Domain{M}^n
  \to \Domain{M}$
  असेल, तर
  $\Subst{s}{f}{u}$
  हे पुढीलप्रमाणे
  परिभाषित केलेले
  चल
  मूल्यांकन आहे:
  \[\Subst{s}{f}{u}(y) = \begin{cases}
    f & \text{जर } y \ident u\\
    s(y) & \text{अन्यथा}.
  \end{cases}\]
  प्रत्येक बाबतीत
  $y$ हे कोणतेही
  प्रथम-क्रम किंवा
  द्वितीय-क्रम
  चल
  असू शकते.
\end{defn}

\begin{defn}[पूर्ति]
द्वितीय-क्रम
सूत्रे~$\varphi$
साठी पूर्तीची
व्याख्या
\ref{fol:syn:sat:defn:satisfaction}
प्रमाणेच आहे;
त्यात पुढील
खंड वाढवले
आहेत:
\begin{enumerate}
\item \indcase{\varphi}{\Atom{X^n}{t_1, \dots, t_n}}{$\Sat{M}{\indfrm}[s]$
  हे $\langle \Value{t_1}{M}[s], \dots, \Value{t_n}{M}[s] \rangle \in
  s(X^n)$ असेल तर आणि तरच खरे.}
\item
  \indcase{\varphi}{\lforall[X][\psi]}{$\Sat{M}{\indfrm}[s]$ हे प्रत्येक $M
  \subseteq \Domain{M}^n$ साठी $\Sat{M}{\psi}[\Subst{s}{M}{X}]$ असेल
  तर आणि तरच खरे.}

\item
  \indcase{\varphi}{\lexists[X][\psi]}{$\Sat{M}{\indfrm}[s]$ हे किमान
  एका $M \subseteq \Domain{M}^n$ साठी
  $\Sat{M}{\psi}[\Subst{s}{M}{X}]$ असेल तर आणि तरच खरे.}

\item
  \indcase{\varphi}{\lforall[u][\psi]}{$\Sat{M}{\indfrm}[s]$ हे प्रत्येक
  $f\colon \Domain{M}^n \to \Domain{M}$ साठी
  $\Sat{M}{\psi}[\Subst{s}{f}{u}]$ असेल तर आणि तरच खरे.}

\item
  \indcase{\varphi}{\lexists[u][\psi]}{$\Sat{M}{\indfrm}[s]$ हे किमान
  एका $f\colon \Domain{M}^n \to \Domain{M}$ साठी
  $\Sat{M}{\psi}[\Subst{s}{f}{u}]$ असेल तर आणि तरच खरे.}
\end{enumerate}
\end{defn}

\begin{ex}
  सूत्र
  $\lforall[z][(\Atom{X}{z} \liff \lnot
    \Atom{Y}{z})]$
  विचारात घ्या.
  त्यात द्वितीय-क्रम
  संख्यक नाहीत,
  पण द्वितीय-क्रम
  चले
  $X$ आणि~$Y$
  आहेत (येथे ती
  एक-स्थानी मानली
  आहेत).
  याच्याशी संबंधित
  प्रथम-क्रम
  वाक्य
  $\lforall[z][(\Atom{P}{z} \liff \lnot \Atom{R}{z})]$
  असे म्हणते की
  आधारक्षेत्रातील प्रत्येक वस्तू
  $P$ आणि $R$ यांच्या
  अर्थनिर्धारणांपैकी नेमक्या एकात येते:
  ती पहिल्यात येते, तर आणि तरच
  ती दुसऱ्यात येत नाही.
  संरचना
  मध्ये
  विधेय~$P$
  चे अर्थनिर्धारण
  $\Assign{P}{M}$
  याने दिलेले
  असते. पण
  $X$ आणि~$Y$
  सारख्या द्वितीय-क्रम
  चले
  चे अर्थनिर्धारण
  संरचना
  स्वतः देत नाही;
  ते चल
  मूल्यांकन देते.
  हे द्वितीय-क्रम
  सूत्र
  वाक्य
  नाही (त्यात
  $X$ आणि~$Y$
  ही मुक्त
  चले
  आहेत);
  म्हणून त्याची
  पूर्ति फक्त
  संरचना~$\Struct{M}$
  आणि त्यासोबत
  चल
  मूल्यांकन~$s$
  यांच्या सापेक्ष
  ठरते.

  $\Sat{M}{\lforall[z][(\Atom{X}{z} \liff \lnot \Atom{Y}{z})]}[s]$
  हे $s(X)$ मधील
  घटक आणि $s(Y)$ मधील
  घटक एकमेकांपासून वेगळे असतील
  आणि दोन्ही संच मिळून संपूर्ण आधारक्षेत्र झाकतील,
  तर आणि तरच खरे असते; म्हणजे
  $s(Y) = \Domain{M} \setminus s(X)$
  असेल तर आणि
  तरच खरे असते.
  उदाहरणार्थ,
  $\Domain{M} = \{1, 2, 3\}$
  घ्या.
  येथे विधेये,
  फलने
  किंवा स्थिरांक
  येत नसल्यामुळे
  $\Struct{M}$
  चे आधारक्षेत्र
  एवढेच संबंधित
  आहे.
  आता
  $s_1(X) = \{1, 2\}$
  आणि
  $s_1(Y) = \{3\}$
  असताना
  $\Sat{M}{\lforall[z][(\Atom{X}{z}
      \liff \lnot \Atom{Y}{z})]}[s_1]$
  खरे आहे.

  याउलट,
  $s_2(X) = \{1, 2\}$
  आणि
  $s_2(Y) = \{2, 3\}$
  असल्यास
  $\Sat/{M}{\lforall[z][(\Atom{X}{z} \liff \lnot
      \Atom{Y}{z})]}[s_2]$.
  कारण
  $\Sat{M}{\Atom{X}{z}}[\Subst{s_2}{2}{z}]$
  (कारण
  $2 \in \Subst{s_2}{2}{z}(X)$),
  पण
  $\Sat/{M}{\lnot \Atom{Y}{z}}[\Subst{s_2}{2}{z}]$
  (कारण
  $2 \in
  \Subst{s_2}{2}{z}(Y)$
  सुद्धा).
\end{ex}

\begin{ex}
  $\Sat{M}{\lexists[Y][(\lexists[y][\Atom{Y}{y}] \land
  \lforall[z][(\Atom{X}{z} \liff \lnot \Atom{Y}{z})])]}[s]$
  हे तेव्हाच खरे
  असते जेव्हा
  असा $N \subseteq \Domain{M}$
  असतो की
  $\Sat{M}{(\lexists[y][\Atom{Y}{y}] \land \lforall[z][(\Atom{X}{z}
  \liff \lnot \Atom{Y}{z})])}[\Subst{s}{N}{Y}]$.
  आणि असे
  $N \neq \emptyset$
  असताना
  ($\Sat{M}{\lexists[y][\Atom{Y}{y}]}[\Subst{s}{N}{Y}]$
  मिळण्यासाठी)
  आणि आधीच्या
  उदाहरणाप्रमाणे
  $N = \Domain{M} \setminus s(X)$
  असताना घडते.
  दुसऱ्या शब्दांत,
  $\Sat{M}{\lexists[Y][(\lexists[y][\Atom{Y}{y}] \land
  \lforall[z][(\Atom{X}{z} \liff \lnot \Atom{Y}{z})])]}[s]$
  हे
  $\Domain{M} \setminus s(X)$
  रिकामे नसेल,
  म्हणजे
  $s(X) \neq
  \Domain{M}$
  असेल, तर आणि
  तरच खरे असते.
  म्हणून
  $\Domain{M}
  = \{1, 2, 3\}$
  आणि
  $s(X) = \{1, 2\}$
  असल्यास
  सूत्र
  ची पूर्ति होते;
  पण
  $s(X) = \{1, 2, 3\}
  = \Domain{M}$
  असल्यास होत नाही.

  $s(X) = \Domain{M}$
  असताना या
  सूत्र
  ची पूर्ति
  होत नसल्यामुळे,
  पुढील
  वाक्य
  \[  \lforall[X][\lexists[Y][(\lexists[y][\Atom{Y}{y}] \land
      \lforall[z][(\Atom{X}{z} \liff \lnot \Atom{Y}{z})])]]
  \]
  कधीच पूर्ण
  होत नाही:
  कोणतीही
  संरचना~$\Struct{M}$
  घेतली तरी
  $s(X) = \Domain{M}$
  हे मूल्यांकन
  त्या
  वाक्य
  ला असत्य
  ठरवेल.
  याउलट,
  पुढील वाक्य
  \[  \lexists[X][\lexists[Y][(\lexists[y][\Atom{Y}{y}] \land
      \lforall[z][(\Atom{X}{z} \liff \lnot \Atom{Y}{z})])]]
  \]
  कोणत्याही
  मूल्यांकन~$s$
  च्या सापेक्ष
  पूर्ण होते,
  कारण
  $M \subseteq \Domain{M}$
  पण
  $M \neq \Domain{M}$
  असा संच नेहमी
  निवडता येतो
  (उदा.,
  $M = \emptyset$).
\end{ex}

\begin{ex}
  द्वितीय-क्रम
  वाक्य~
  $\lforall[X][\lforall[y][X(y)]]$
  असे म्हणते की
  प्रत्येक
  $1$-स्थानी
  संबंध, म्हणजे
  प्रत्येक
  गुणधर्म, प्रत्येक
  वस्तूला लागू
  होतो.
  हे स्पष्टपणे
  कधीच खरे नसते:
  प्रत्येक~$\Struct{M}$
  मध्ये
  $s(X) = \emptyset$
  आणि
  $s(y) = a \in
  \Domain{M}$
  असे चर-मूल्यांकन~$s$
  घेतल्यास
  $\Sat/{M}{X(y)}[s]$
  मिळते.
  याचा अर्थ
  $\varphi \lif
  \lforall[X][\lforall[y][X(y)]]$
  हे द्वितीय-क्रम
  तर्कशास्त्रात
  $\lnot \varphi$
  शी सममूल्य
  आहे; म्हणजे
  $\Sat{M}{\varphi \lif
  \lforall[X][\lforall[y][X(y)]]}$
  हे
  $\Sat{M}{\lnot \varphi}$
  असेल तर आणि
  तरच खरे.
  दुसऱ्या शब्दांत,
  द्वितीय-क्रम
  तर्कशास्त्रात
  $\lforall$
  आणि~$\lif$
  वापरून
  $\lnot$
  परिभाषित
  करता येते.
\end{ex}

\begin{prob}
  द्वितीय-क्रम
  तर्कशास्त्रात
  $\lforall$
  आणि~$\lif$
  वापरून इतर
  संयोजके
  परिभाषित
  करता येतात
  हे दाखवा:
  \begin{enumerate}
    \item येथे X हे संयोजित करावयाच्या दोन्ही सूत्रांमध्ये मुक्त नसलेले
    एकस्थानी विधेय चल निवडा. द्वितीय-क्रम तर्कशास्त्रात $\varphi \land \psi$ हे
    $\lforall[X][(\varphi \lif (\psi \lif \lforall[x][X(x)])\lif
    \lforall[x][X(x)])]$ शी सममूल्य आहे हे सिद्ध करा.
    \item फक्त $\lforall$ आणि $\lif$ वापरून
    $\varphi \lor \psi$ शी सममूल्य द्वितीय-क्रम सूत्र शोधा.
  \end{enumerate}
\end{prob}











\section{चिन्हार्थविषयक संकल्पना}\label{sol:syn:sem:sec}

\begin{explain}
\emph{वैधता},
\emph{तार्किक निष्पन्नता}
आणि
\emph{पूर्ततायोग्यता}
या मध्यवर्ती
तार्किक संकल्पना
द्वितीय-क्रम
तर्कशास्त्रासाठी
प्रथम-क्रम
तर्कशास्त्राप्रमाणेच
परिभाषित होतात;
फरक इतकाच की
आधारभूत पूर्ति-संबंध
आता द्वितीय-क्रम
सूत्रे
साठीचा असतो.
अर्थात, द्वितीय-क्रम
वाक्य
हे असे
सूत्र
असते ज्यातील
विधेय-चल आणि
फलन-चल यांसह
सर्व चले
बद्ध असतात.
\end{explain}

\begin{defn}[वैधता]
वाक्य $\varphi$
\emph{वैध},
$\Entails \varphi$,
असते जर आणि
फक्त जर प्रत्येक
संरचना~$\Struct M$
साठी $\Sat{M}{\varphi}$.
\end{defn}

\begin{defn}[तार्किक निष्पन्नता]
वाक्यांच्या
संच~$\Gamma$
मधून वाक्य~$\varphi$
\emph{तार्किकरीत्या
निष्पन्न होते},
$\Gamma
\Entails \varphi$,
जर आणि फक्त जर
$\Sat{M}{\Gamma}$
असलेल्या प्रत्येक
संरचना~$\Struct M$
साठी
$\Sat{M}{\varphi}$.
\end{defn}

\begin{defn}[पूर्ततायोग्यता]
एखाद्या
संरचना~$\Struct M$
साठी
$\Sat{M}{\Gamma}$
असेल, तर वाक्यांचा
संच~$\Gamma$
\emph{पूर्ततायोग्य}
असतो.
$\Gamma$
पूर्ततायोग्य नसेल,
तर त्याला
\emph{अपूर्ततायोग्य}
म्हणतात.
\end{defn}











\section{अभिव्यक्तीचे सामर्थ्य}\label{sol:syn:exp:sec}

\begin{explain}
द्वितीय-क्रम
चलांवरील
संख्यकीकरणामुळे
भाषेचे अभिव्यक्तीचे
सामर्थ्य प्रथम-क्रम
तर्कशास्त्राच्या
तुलनेत प्रचंड
वाढते.
द्वितीय-क्रम
अस्तित्ववाचक
संख्यकीकरणाने
विशिष्ट गुणधर्म
असलेली फलने
किंवा संबंध
अस्तित्वात आहेत,
असे म्हणता येते.
प्रथम-क्रम
तर्कशास्त्रात
हे करण्याचा
एकमेव मार्ग
म्हणजे त्यासाठी
अतार्किक चिन्ह
(म्हणजे फलन
किंवा विधेय)
निर्दिष्ट करणे.
द्वितीय-क्रम
वैश्विक संख्यकीकरणाने
आधारक्षेत्राचे
सर्व उपसंच,
त्यावरील सर्व
संबंध किंवा
प्रांत वरून
प्रांत मध्ये
जाणारी सर्व
फलने यांना
एखादा गुणधर्म
आहे, असे
म्हणता येते.
प्रथम-क्रम
तर्कशास्त्रात
भाषेतील अतार्किक
चिन्हांपैकी
एखाद्याला
नेमून दिलेल्या
उपसंचांना,
संबंधांना किंवा
फलनांनाच
गुणधर्म आहे,
एवढेच म्हणता
येते.
एखादा गुणधर्म
असलेले उपसंच,
संबंध किंवा
फलने अस्तित्वात
आहेत किंवा
त्या सर्वांना
तो गुणधर्म
आहे असे म्हणताना,
त्या गुणधर्माचे
वर्णन करतानाही
द्वितीय-क्रम
संख्यकीकरण
वापरता येते.
यामुळे प्रथम-क्रम
तर्कशास्त्रात
परिभाषित न
होणारे संबंध
परिभाषित करता
येतात आणि
तेथे व्यक्त न
होणारे आधारक्षेत्राचे
गुणधर्म व्यक्त
करता येतात.
\end{explain}

\begin{defn}
$\Struct{M}$ ही
भाषा~$\Lang{L}$
साठी संरचना
असेल, तर
$R \subseteq \Domain{M}^2$
हा संबंध~$\Lang{L}$
मध्ये \emph{परिभाष्य}
असतो, जर फक्त
$x$ आणि~$y$
हीच मुक्त चले
असलेले असे
सूत्र~$\varphi_R(x, y)$
असते की
$s(x) = a$
आणि $s(y) = b$
असताना
$R(a, b)$
खरे असते
(म्हणजे
$\tuple{a,b} \in R$),
जर आणि फक्त
जर
$\Sat{M}{\varphi_R(x, y)}[s]$.
\end{defn}

\begin{ex}
प्रथम-क्रम
तर्कशास्त्रात
एकरूपता संबंध~
$\Id{\Domain{M}}$
(म्हणजे
$\Setabs{\tuple{a,a}}{a \in
  \Domain{M}}$)
सूत्र~$\eq[x][y]$
याने परिभाषित
करता येतो.
द्वितीय-क्रम
तर्कशास्त्रात
हा संबंध
\emph{$\eq$ शिवाय}
परिभाषित
करता येतो.
कारण $a$
आणि $b$
हे~$\Domain{M}$
चे तेच
घटक
असतील, तर
ते~$\Domain{M}$
च्या त्याच
उपसंचांचे
घटक
असतात (कारण
संच त्यांतील
घटक
नी ठरतात).
उलट $a$
आणि $b$
भिन्न असतील,
तर ते त्याच
उपसंचांचे
घटक
नसतात:
उदा.,
$a \in \{a\}$
पण $a \neq b$
असताना
$b
\notin \{a\}$.
म्हणून
“~$\Domain{M}$
च्या त्याच
उपसंचांचे
घटक
असणे” हा
संबंध $a$
आणि $b$
यांच्याबाबत
$a = b$
असेल तर आणि
तरच लागू
होतो.
हा संबंध
द्वितीय-क्रम
तर्कशास्त्रात
व्यक्त करता
येतो, कारण
~$\Domain{M}$
च्या सर्व
उपसंचांवर
संख्यकीकरण
करता येते.
म्हणून पुढील
सूत्र
$\Id{\Domain{M}}$
परिभाषित करते:
\[\lforall[X][(X(x) \liff X(y))]
\]
\end{ex}

\begin{prob}
$\lforall[X][(X(x) \lif X(y))]$
(लक्षात घ्या:
$\liff$ नव्हे,
$\lif$!)
हे सूत्र
$\Id{\Domain{M}}$
परिभाषित करते,
हे दाखवा.
\end{prob}

\begin{ex}
$R$ हे
द्विस्थानी
विधेय
असेल, तर
$\Assign{R}{M}$
हा~$\Domain{M}$
वरील द्विस्थानी
संबंध असतो.
थोडे गोंधळात
टाकणारे असले,
तरी~$R$ हेच
चिन्ह त्या
विधेय
साठी~$R$ आणि
संबंध~$\Assign{R}{M}$
साठीही वापरू.
$R$ चे
\emph{संक्रमक
संवरण}~$R^+$
हा $a$
आणि $b$
यांच्यात तेव्हाच
लागू होणारा
संबंध आहे
जेव्हा काही
$c_1$, \dots, $c_k$
साठी
$R(a,c_1)$,
$R(c_1, c_2)$,
\dots, $R(c_k,b)$
लागू होतात.
यात $k = 0$
हा प्रसंगही
येतो, म्हणजे
$R(a,b)$
लागू असेल,
तर $R^+(a,b)$
ही लागू असतो.
याचा अर्थ
$R \subseteq R^+$.
खरे तर
$R^+$ हा~$R$
समाविष्ट करणारा
सर्वांत लहान
संक्रमक संबंध
आहे.
द्वितीय-क्रम
तर्कशास्त्रात
$X$ हा~$R$
समाविष्ट करणारा
संक्रमक संबंध
आहे असे
म्हणता येते:
\begin{multline*}
  \psi_R(X) \ident \lforall[x][\lforall[y][(R(x,y) \lif X(x, y))]] \land {}\\
\lforall[x][\lforall[y][\lforall[z][((X(x,y) \land X(y,z)) \lif X(x,
      z))]]].
\end{multline*}
पहिला संयुक्तांश
$R \subseteq X$
सांगतो आणि
दुसरा $X$
संक्रमक आहे
असे सांगतो.

$X$ हा
असा सर्वांत
लहान संबंध
आहे असे
म्हणणे म्हणजे
$R$ समाविष्ट
करणाऱ्या
आणि संक्रमक
असलेल्या
प्रत्येक संबंधात
तो समाविष्ट
आहे असे
म्हणणे.
म्हणून
पुढील
सूत्र
याने~$R$
चे संक्रमक
संवरण
परिभाषित
करता येते:
\[R^+(X) \ident \psi_R(X) \land \lforall[Y][(\psi_R(Y) \lif
  \lforall[x][\lforall[y][(X(x, y) \lif Y(x,y))]])].
\]
$\Sat{M}{R^+(X)}[s]$
हे
$s(X) = R^+$
असेल तर आणि
तरच खरे.
$R$ चे
संक्रमक संवरण
प्रथम-क्रम
तर्कशास्त्रात
व्यक्त करता
येत नाही.
\end{ex}












\section{अनंत आणि गणनीय
  प्रांत यांचे वर्णन}\label{sol:syn:siz:sec}


येथे निवडीचे स्वयंसिद्धक मानणारा नेहमीचा
संचसैद्धान्तिक संदर्भ घेतला आहे. त्यात डेडेकिंड अनंतता
आणि सर्वसाधारण अनंतता सममूल्य असतात.
संच~$M$ हा
(डेडेकिंडच्या अर्थाने)
अनंत असतो,
जर आणि फक्त जर
$f\colon M
\to M$ असे
एकास-एक
फलन असते जे
आच्छादक
नसते; म्हणजे
$\ran{f} \neq M$.
प्रथम-क्रम
तर्कशास्त्रात
एक-स्थानी
फलन~$f$
घेऊन
संरचना~$\Struct{M}$
मध्ये त्याला
नेमलेले फलन
$\Assign{f}{M}$
हे एकास-एक
आहे आणि
$\ran{f} \neq
\Domain{M}$
आहे असे
म्हणता येते:
\[\lforall[x][\lforall[y][(\eq[f(x)][f(y)] \to \eq[x][y])]] \land
\lexists[y][\lforall[x][\eq/[y][f(x)]]].
\]
$\Struct{M}$
ही या
वाक्य
ची पूर्ति करत
असेल, तर
$\Assign{f}{M}:
\Domain{M} \to \Domain{M}$
हे एकास-एक
असते, म्हणून
$\Domain{M}$
अनंत असलेच
पाहिजे.
$\Domain{M}$
अनंत असेल
आणि म्हणून
असे फलन
अस्तित्वात असेल,
तर ते फलन
$\Assign{f}{M}$
म्हणून घेता
येते आणि
$\Struct{M}$
ही त्या
वाक्य
ची पूर्ति करेल.
मात्र यासाठी
आपल्या भाषेत
या हेतूने
वापरण्यात येणारे
अतार्किक चिन्ह~$f$
असणे आवश्यक
आहे.
द्वितीय-क्रम
तर्कशास्त्रात
असे फलन
\emph{अस्तित्वात आहे}
असे थेट
म्हणता येते.
यासाठी
आता~$f$
लागत नाही,
आणि शुद्ध
द्वितीय-क्रम
तर्कशास्त्रातील
पुढील
वाक्य
मिळते:
\[\fn{Inf} \ident \lexists[u][(\lforall[x][\lforall[y][(\eq[u(x)][u(y)]
      \lif \eq[x][y])]] \land \lexists[y][\lforall[x][\eq/[y][u(x)]]])].
\]
$\Sat{M}{\fn{Inf}}$
हे
$\Domain{M}$
अनंत असेल
तर आणि तरच
खरे असते.
मग
$\fn{Fin} \ident \lnot \fn{Inf}$
असे परिभाषित
करता येते;
$\Sat{M}{\fn{Fin}}$
हे
$\Domain{M}$
सांत असेल
तर आणि तरच
खरे असते.
शुद्ध प्रथम-क्रम
तर्कशास्त्रातील
कोणतेही एक
वाक्य
प्रांत
अनंत आहे
हे व्यक्त
करू शकत
नाही, मात्र
अशा वाक्यांचा
अनंत संच
ते करू शकतो.
शुद्ध प्रथम-क्रम
तर्कशास्त्रातील
वाक्ये
चा असा
कोणताही संच
नाही की
त्याची
संरचना
मध्ये पूर्ति
होते जर आणि
फक्त जर तिचे
आधारक्षेत्र
सांत आहे.

\begin{prop}
$\Sat{M}{\fn{Inf}}$
हे
$\Domain{M}$
अनंत असेल
तर आणि तरच
खरे असते.
\end{prop}

\begin{proof}
$\Sat{M}{\fn{Inf}}$
असेल तर आणि
तरच काही~$s$
साठी
  $\Sat{M}{\lforall[x][\lforall[y][(\eq[u(x)][u(y)] \lif \eq[x][y])]]
    \land \lexists[y][\lforall[x][\eq/[y][u(x)]]]}[s]$.
  असे असेल,
  तर $s(u)$
  हे एकास-एक
  फलन आहे
  आणि काही
  $y
  \in \Domain{M}$
  हे~$s(u)$
  च्या व्याप्तीत
  नाही.
  उलट
  $f\colon \Domain{M} \to \Domain{M}$
  असे
  एकास-एक
  फलन असेल
  की
  $\ran{f}
  \neq \Domain{M}$,
  तर
  $s(u) = f$
  असे चर-मूल्यांकन
  मिळते.
\end{proof}

संच $M$
गणनीय
असतो, जर
त्यातील
घटक
ची पुनरावृत्ती
नसलेली
(पण शक्यतो
सांत असलेली)
पुढील
प्रगणना असते:
\[m_0, m_1, m_2, \dots
\]
अशी प्रगणना
असते जर
आणि फक्त
जर
घटक
$z \in M$
आणि
$f\colon M \to M$
असे फलन
असते की
$z$, $f(z)$,
$f(f(z))$, \dots,
हे~$M$
चे सर्व
घटक
असतात.
कारण प्रगणना
असेल, तर
$z = m_0$
आणि
$f(m_k) = m_{k+1}$
(किंवा
$m_k$ हा
प्रगणनेतील
शेवटचा
घटक
असेल, तर
$f(m_k) = m_k$)
हे अपेक्षित
घटक
आणि फलन
ठरतात.
उलट असे
$z$ आणि
$f$ असतील,
तर
$z$, $f(z)$,
$f(f(z))$, \dots,
या यादीतील पुनरावृत्ती
काढल्यावर~$M$ ची प्रगणना
मिळते आणि
$M$ हा
गणनीय
असतो.
$z$ आणि~$f$
यांचे अस्तित्व
द्वितीय-क्रम
तर्कशास्त्रात
व्यक्त करून
असे
वाक्य
तयार करता
येते की
ते
संरचना
मध्ये खरे
असते जर
आणि फक्त
जर ती
संरचना
गणनीय
असते:
\[\fn{Count} \ident
\lexists[z][\lexists[u][\lforall[X][((X(z) \land
      \lforall[x][(X(x) \lif X(u(x)))]) \lif \lforall[x][X(x)])]]]
\]

\begin{prop}
$\Sat{M}{\fn{Count}}$
हे
$\Domain{M}$
गणनीय
असेल तर
आणि तरच
खरे असते.
\end{prop}

\begin{proof}
$\Domain{M}$
गणनीय
आहे असे
समजू आणि
$m_0$, $m_1$,
\dots, ही
प्रगणना घेऊ.
पुनरावृत्ती
काढून टाकल्यास
कोणतेही
$m_k$ दोनदा
येणार नाही
याची खात्री
करता येते.
$f(m_k) = m_{k+1}$
असे परिभाषित
करू (सांत
प्रगणनेच्या
शेवटच्या घटकावर
फलन तोच घटक
परत देते)
आणि
$s(z) = m_0$,
$s(u) = f$
घेऊ.
आता पुढील
गोष्ट दाखवू:
\[\Sat{M}{\lforall[X][((X(z) \land \lforall[x][(X(x) \lif X(u(x)))])
    \lif \lforall[x][X(x)])]}[s]
\]
$M \subseteq \Domain{M}$
हा कोणताही
संच असू
द्या.
तसेच
$\Sat{M}{(X(z) \land \lforall[x][(X(x) \lif
X(u(x)))])}[\Subst{s}{M}{X}]$
असे धरू.
मग
$\Subst{s}{M}{X}(z) \in M$
आणि
$x \in M$
असेल तेव्हा
$(\Subst{s}{M}{X}(u))(x) \in M$
ही.
दुसऱ्या
शब्दांत,
$\varAssign{\Subst{s}{M}{X}}{s}{X}$
असल्यामुळे
$m_0 \in M$
आणि
$x \in M$
असेल तर
$f(x) \in M$;
म्हणून
$m_0 \in M$,
$m_1 = f(m_0) \in M$,
$m_2 = f(f(m_0)) \in M$
इत्यादी.
त्यामुळे
$M = \Domain{M}$
आणि म्हणून
$\Sat{M}{\lforall[x][X(x)]}[\Subst{s}{M}{X}]$.
$M \subseteq
\Domain{M}$
कोणताही
असल्यामुळे
सिद्धता पूर्ण:
$\Sat{M}{\fn{Count}}$.

आता
$\Sat{M}{\fn{Count}}$,
म्हणजे
\[\Sat{M}{\lforall[X][((X(z) \land \lforall[x][(X(x) \lif X(u(x)))])
    \lif \lforall[x][X(x)])]}[s]
\]
काही~$s$
साठी आहे
असे समजू.
$m = s(z)$
आणि
$f = s(u)$
घेऊ आणि
$M = \{m,
f(m), f(f(m)), \dots\}$
हा संच
विचारात
घेऊ.
असा
परिभाषित
केलेला
$M$
स्पष्टपणे
गणनीय
आहे.
मग
\[\Sat{M}{(X(z) \land \lforall[x][(X(x) \lif X(u(x)))])
    \lif \lforall[x][X(x)]}[\Subst{s}{M}{X}]
\]
हे गृहीतकावरून
मिळते.
तसेच
$\Sat{M}{X(z)}[\Subst{s}{M}{X}]$,
कारण
$M \ni m =
\Subst{s}{M}{X}(z)$;
आणि
$\Sat{M}{\lforall[x][(X(x) \lif
X(u(x)))]}[\Subst{s}{M}{X}]$
सुद्धा,
कारण
$x \in M$
असेल तेव्हा
$f(x) \in
M$.
म्हणून
पूर्वपक्ष
आणि शर्त
दोन्हींची
पूर्ति
झाल्याने
उत्तरपक्षाचीही
पूर्ति झाली
पाहिजे:
$\Sat{M}{\lforall[x][X(x)]}[\Subst{s}{M}{X}]$.
पण याचा
अर्थ
$M = \Domain{M}$.
म्हणून
$\Domain{M}$
गणनीय
आहे, कारण
$M$ व्याख्येने
तसाच आहे.
\end{proof}

\begin{prob}
वाक्य
$\fn{Inf} \land \fn{Count}$
हे नेमक्या
गणनीय अनंत
आधारक्षेत्रांत
खरे असते.
$\fn{Count}$
ची व्याख्या
बदलून
आधारक्षेत्र
गणनीय अनंत
आहे हे
थेट व्यक्त
करणारे
वेगळे
वाक्य
मिळवा आणि
ते तसे
करते हे
सिद्ध करा.
\end{prob}



